\subsection{False Positives and False Negatives}
We distinguish the predicate-selected set from the fixed-budget set supplied to the downstream agent. Let $C_q$ be the retrieved candidate set, $P_q$ the predicate-selected set, and $G_q$ the annotated required set. In this error analysis we use the authors' nonstandard error labels: a false positive (FP) is a required skill whose evaluated predicate is false, and a false negative (FN) is a nonrequired skill whose evaluated predicate is true. Thus predicate FP counts $|(G_q\cap C_q)\setminus P_q|$ and predicate FN counts $|P_q\setminus G_q|$. We retain standard relevance-based TP as $|P_q\cap G_q|$ and TN as $|C_q\setminus(P_q\cup G_q)|$; this hybrid naming is not a standard confusion-matrix convention.

For set-level accounting over $C_q\cup G_q$, FP counts all missing required skills, $|G_q\setminus P_q|$, including retrieval misses $|G_q\setminus C_q|$ listed separately below. A retrieval miss has no evaluated predicate and is not a predicate rejection. Unretrieved non-gold catalog skills do not inflate TN. For a supplied top-five set, the same set-level error labels describe missing required and extra nonrequired skills, not necessarily the underlying predicate decisions.

\begin{table*}[t]
\centering
\small
\begin{tabular}{llrrrrrr}
\toprule
Dataset & Retriever & Tasks & TP & TN & FN & FP & Retrieval miss \\
\midrule
TheoremQA & BM25 & 105 & 61 & 5104 & 49 & 44 & 8 \\
TheoremQA & Vector & 105 & 63 & 5069 & 79 & 42 & 3 \\
TheoremQA & SkillSight & 105 & 66 & 5069 & 78 & 39 & 2 \\
TheoremQA & Pooled & 315 & 190 & 15242 & 206 & 125 & 13 \\
MedCalc-Bench & BM25 & 165 & 135 & 8085 & 7 & 30 & 7 \\
MedCalc-Bench & Vector & 165 & 138 & 8077 & 8 & 27 & 0 \\
MedCalc-Bench & SkillSight & 165 & 139 & 8079 & 6 & 26 & 0 \\
MedCalc-Bench & Pooled & 495 & 412 & 24241 & 21 & 83 & 7 \\
\bottomrule
\end{tabular}
\caption{Final predicate-selected-set counts with the authors' nonstandard FP/FN labels and standard relevance-based TP/TN. Pooled rows sum task--retriever observations, not distinct tasks or independent training repeats. Retrieval misses are the subset of set-level FP absent from the top-50 candidate pool, not false predicate judgments.}
\label{tab:app-error-confusion}
\end{table*}

Each retriever is evaluated on 105 TheoremQA tasks and 165 MedCalc-Bench tasks, with one annotated required skill per task. Table~\ref{tab:app-error-confusion} shows different error directions: false negatives (extra nonrequired skills) account for 62.24\% of TheoremQA's FP+FN occurrences, whereas set-level false positives (missing required skills) account for 79.81\% of MedCalc-Bench's. Of their respective set-level false positives, 112 and 76 concern retrieved skills whose predicates did not activate; the remaining 13 and 7 are retrieval misses. Failure to activate is an outcome, not by itself evidence that a predicate is too narrow.

For TheoremQA, we additionally review each error against the task, predicate, and skill scope (Table~\ref{tab:app-error-attribution}). Incorrect judgment means that the saved Boolean conflicts with the predicate's meaning for the task. A broad predicate admits work outside the skill's scope; a restrictive or misdirected predicate excludes an annotated required use. Annotation/scope mismatch covers plausible applicability beyond the annotated required set, and uncertain cases remain ambiguous. These categories are retrospective semantic assessments rather than explanations generated by the checker. They separate checker mistakes from predicate-design and annotation issues instead of assigning all errors to one mechanism.

\begin{table*}[t]
\centering
\small
\begin{tabular}{lrrrr}
\toprule
Category & FN & FP & Total & Share (\%) \\
\midrule
Incorrect predicate judgment & 83 & 31 & 114 & 34.4 \\
Overly broad predicate & 14 & 0 & 14 & 4.2 \\
Restrictive or misdirected predicate & 0 & 70 & 70 & 21.1 \\
Plausible annotation/scope mismatch & 76 & 2 & 78 & 23.6 \\
Ambiguous attribution & 33 & 9 & 42 & 12.7 \\
Retrieval miss & 0 & 13 & 13 & 3.9 \\
\bottomrule
\end{tabular}
\caption{Retrospective semantic categorization of all 331 TheoremQA task--skill selection errors across the three retrievers. Shares use all 331 occurrences as denominator; examples alone are not used to estimate frequencies.}
\label{tab:app-error-attribution}
\end{table*}
